% ======================================================================================================
\section{Beam-loaded plant and recorded response}\label{app:derivation}
% ======================================================================================================

With the cavity voltages $\vect v_k=\mat Z_k(\vect i_g-I_{b1}\hat{\vect b}\,\Delta\phi)$ and Eq.~\eqref{eq:beam},
\begin{equation}
\Bigl(s^2+2\arad s+\omega_s^2\Bigr)\Delta\phi=\kappa_1\hat{\vect p}^\mathsf{T}\mat Z_\Sigma\vect i_g-\kappa_1I_{b1}\bigl(\hat{\vect p}^\mathsf{T}\mat Z_\Sigma\hat{\vect b}\bigr)\Delta\phi ,
\end{equation}
so $\Delta\phi=\kappa_1\hat{\vect p}^\mathsf{T}\mat Z_\Sigma\vect i_g/D(s)$ and $\vect v_m=\mat Z_m\vect i_g-I_{b1}\mat Z_m\hat{\vect b}\,\Delta\phi$ gives Eq.~\eqref{eq:plant}. The term $\kappa_1I_{b1}\hat{\vect p}^\mathsf{T}\mat Z_\Sigma\hat{\vect b}$ equals $-j\kappa[Z_\Sigma(s)-\bar Z_\Sigma(s)]$ with $\kappa=\kappa_1I_{b1}/2$, so that $D(s)$ is the open-loop characteristic function of the beam-loaded cavity~\cite{Pedersen1975}; multiplied by the cavity denominator it is Pedersen's quartic.

The loop of Eq.~\eqref{eq:loop} places the klystron between the drive chain and the cavities. Blocks with a real impulse response commute with every $2\times2$ matrix, so only the order of the rotation $\mat R(\theta)$, the plant $\mat P_b$ and the klystron $\mat K$ matters. The noise is injected into the drive in the RFP frame; with $\psi$ the phase from the RFP drive frame to the cavity frame, conjugating every signal by $\mat R(\psi)$ moves the loop to the cavity frame and turns the injection angle $\alpha_\mathrm{RFP}$ into $\alpha=\alpha_\mathrm{RFP}+\psi$. The loop phase $\theta$ includes $\psi$ and the controller rotation; a change of the loop-phase setting changes $\theta$ but not $\alpha$, while a change of the klystron's insertion phase with its cathode voltage changes both.

If the excitation is one-dimensional, with a real amplitude $r(t)$ along a fixed angle $\alpha$, $\vect n(t)=r(t)(\cos\alpha,\sin\alpha)^\mathsf{T}$, the complex input is $n=e^{j\alpha}r$ and $\overline{n(-f)}=e^{-2j\alpha}n(f)$; the widely linear output $x=H_1n+H_2\overline{n(-f)}$ is then $(H_1+H_2e^{-2j\alpha})\,n$, Eq.~\eqref{eq:recorded}. The recorded response is fully coherent with the reference, but it is one projection: $H_1$ and $H_2$ cannot be separated from one record, and parameters that act only through the image response $H_2$---the klystron's asymmetry above all---are confounded with the loop parameters (Sec.~\ref{sec:ident}).

The closed-loop poles are the zeros of $\det[\mat 1+\mat M(s)]D(s)$. With $\mat F=\mat K\,Gc\,\mat R(\theta)/R_L$, Sylvester's identity gives $\det(\mat 1+\mat M)=\det(\mat 1+\mat F\mat P_b)$, and since $\mat P_b$ is $\mat Z_m$ minus a rank-one term, the matrix determinant lemma gives the exact factorization
\begin{equation}
\det[\mat 1+\mat M(s)]\,D(s)=\det(\mat 1+\mat F\mat Z_m)\,D_0(s),
\end{equation}
where $D_0(s)=s^2+2\arad s+\omega_s^2+\kappa_1I_{b1}\hat{\vect p}^\mathsf{T}\mat Z_b\hat{\vect b}$ is the characteristic function of the beam with the closed-loop impedance $\mat Z_b=\mat Z_\Sigma(\mat 1+\mat F\mat Z_m)^{-1}$ that feedback presents to it~\cite{Boussard1985,Wilson1993}, and $\det(\mat 1+\mat F\mat Z_m)$ that of the RF loop without beam. The coupled mode is tracked by Newton continuation from $\Ib=0$, where it is the incoherent oscillator $s=-\arad\pm j(\omega_s^2-\arad^2)^{1/2}$.

% ======================================================================================================
\section{Estimator statistics}\label{app:estimator}
% ======================================================================================================

With the bin offset $r=-n,\dots,n$ ($n=4$) around a centre bin $k$, the LPM fits
\begin{equation}
Y_{k+r}=\hat G\,X_{k+r}+\sum_{q=0}^{2}t_q\,(r/n)^q+E_r
\end{equation}
by least squares~\cite{Pintelon2010}, i.e.\ $\vect Y=\mat K\vect\vartheta+\vect E$ with $\vect\vartheta=(\hat G,t_0,t_1,t_2)^\mathsf{T}$ and $\hat G=\vect e_0^\mathsf{T}(\mat K^\mathsf{H}\mat K)^{-1}\mat K^\mathsf{H}\vect Y$. The neighbourhoods tile the spectrum outward from the carrier; in the one centred on it, the carrier bin, which holds the record's mean, is left out of the fit. Any quadratic transient is annihilated, and with Eq.~\eqref{eq:lpm} the expectation of the estimate is
\begin{align}
\langle\hat G\rangle&=\sum_r c_r\,G_{k+r},\\
\vect c^\mathsf{T}&=\vect e_0^\mathsf{T}(\mat K^\mathsf{H}\mat K)^{-1}\mat K^\mathsf{H}\,\mathrm{diag}(X_{k+r}),
\end{align}
with weights that sum to one; for output noise of variance $\sigma_V^2$ per bin its variance is $\sigma_V^2[(\mat K^\mathsf{H}\mat K)^{-1}]_{00}$. The noise variance is estimated from the residual of the same fit with $G$ quadratic as well (six unknowns, three degrees of freedom), $s_k^2=\sum_r|E_r|^2/3$, averaged over \SI{\lpmSmooth}{kHz}. Holding $G$ constant in the estimate costs no variance and causes no bias, because the model is averaged with the same weights $c_r$; this matters near the coupled mode, where the response changes within one neighbourhood. In the loop band, where several estimates form one data point, the model is evaluated at their weighted mean frequency.

The predicted variance is tested without a model on the scatter of adjacent estimates about a linear trend. Above \SI{100}{kHz}, in groups of eight (\SI{2.8}{kHz}), the weighted residual is \calibMedian{} times its expectation (median over the records; \calibLo--\calibHi). In the near band, in groups of four, the ratio lies between \calibNearLo{} and \calibNearHi: interference lines near $\pm22$ and $\pm\SI{27}{kHz}$ raise the averaged noise level of the estimates around them, which therefore weigh less, and below \SI{5}{kHz} the response itself curves within a group. Within \SI{3}{kHz} of the carrier, where the noise level changes fast and the harmonics of the high-voltage ripple (multiples of \SI{720}{Hz}) stand out, the averaged level is too low on the estimates that contain a harmonic and too high on the others: there the residuals of the selected model are \nearCalLo{} to \nearCalHi{} times their predicted standard deviation (median over the records). Averaging the noise level over fewer bins does not remove this and makes the level itself noisier. The Welch estimate that the routine uses has the right variance per bin, but adjacent bins of overlapping segments are correlated and its resolution is set by the segment length; treating its bins as independent overstates the information they carry, and near the coupled mode the model has to be convolved with the window. The LPM uses the whole record, removes the transient and gives independent estimates.

The model-free estimate of the coupled mode fits, on both sides of the carrier at once (3 to \SI{16}{kHz}), a smooth background (complex quadratic in $f$) plus the residues of one complex pole pair $s=-\sigma\pm j\omega_c$, each averaged with the weights $c_r$ of the estimates.

% ======================================================================================================
\section{Trip records}\label{app:trips}
% ======================================================================================================

At a trip the station's fast acquisition memories freeze at the fault and are written to files. Each of the \tripN{} records of 2021, all of the present tube, holds \num{51200} samples per signal at \SI{10.24}{MS/s}, the \SI{5.0}{ms} before the fault. The signals are the klystron drive and one cavity probe, both as I/Q, and the klystron output, whose amplitude $(I^2+Q^2)^{1/2}$ is linear in the RF amplitude. As for the closed-loop records, the stored I/Q are the complex conjugate of the physical envelope: conjugated, the drive-to-probe response places the cavities below $\fRF$ in every record, its phase-path response being \tripAsymLo--\tripAsymHi{} times larger at \SIrange{-45}{-20}{kHz} than at \SIrange{20}{45}{kHz}. In \tripUsed{} records every signal holds RF; in the others some or all of them are empty. The analysis uses the normal operation up to \SI{0.5}{ms} before the end of each record, earlier if any signal leaves its level first.

The compression is the cross-spectrum of the relative output amplitude with the relative drive amplitude over the drive's auto-spectrum ($H_1$; 2048-point Hann segments with \SI{50}{\percent} overlap), with the ratio of the output's auto-spectrum to the cross-spectrum ($H_2$) as the bound from the other side; the two acquisition channels differ in timing by \SI{\tripDelay}{ns}, which only turns the phase of the ratio. For the AM/PM conversion the relative fluctuation $c$ of the probe is regressed, bin by bin (1024-point segments), on the relative drive amplitude $a_d$ and the drive phase $\varphi_d$: $c=A\,a_d+B\,\varphi_d$. For a memoryless klystron $c=P[(\gA+j\kAP)\,a_d+j\varphi_d]$ for any complex-linear plant $P$, so that $A/B=\kAP-j\gA$ whatever the cavities and their detuning; the beam acts on the probe only near $\pm f_s$, outside the band of \SIrange{60}{250}{kHz} used. Feedback correlates the drive with disturbances of the probe and pulls the ratio towards that of a linear klystron ($A/B\to-j$); the compression the regression returns, \tripGARatio{} against \tripGA{} measured directly, gauges that pull.
